\section{Radial sweeps and finite cylinders}
\label{sec:introduction}

A cylinder in $\R^3$ is a set $B+\R u$, where $u$ is a unit vector
and $B\subset u^\perp$ is a measurable planar set of finite area.
Its cost is $\area(B)$. In a finite cover, the costs are added with
multiplicity. For a convex body $K$, write
\[
 A_{\min}(K)=\min_{\norm{u}=1}\area(\pi_{u^\perp}K).
\]
The half-area question asks whether every finite cylinder cover of $K$
has cost at least $A_{\min}(K)/2$. A regular tetrahedron has a simple
cover by two opposite-edge cylinders attaining this value.
A second covering quantity is the directionwise normalized cost
\[
 \sum_i\frac{\area(B_i)}{\area(\pi_{u_i^\perp}K)},
\]
which divides each base area by the shadow in its own axis direction.

These questions extend plank covering, where the cost of a covering
slab is its width rather than the area of a planar base. Bezdek's
Problem~3.1 attributes the half-area question and the tetrahedral
equality example to Bang~\cite{Bezdek2009}.
Bezdek and Litvak prove that the directionwise normalized cost is at
least \(1/3\) for every three-dimensional convex body and at least \(1\)
for ellipsoids~\cite[Theorem~3.1]{BL2009}. In dimension three, Bezdek
and Khan's \emph{1-Codimensional Cylinder Covering Conjecture} proposes
the lower bound \(1/2\) for all convex bodies
\cite[Definition~3 and Conjecture~4.13]{BK2016}.
Such a normalized bound would imply the half-minimum bound, since
each directional projection area is at least \(A_{\min}(K)\).
The half-area question also appears in Verreault's survey
\cite[Question~4.14]{Verreault2026}. Our construction lowers the
tetrahedron's two-cylinder cost strictly by varying the directions,
while keeping coverage after finite approximation.

Our approximation problem begins before the cylinders are selected.
Let $P$ be a two-dimensional Euclidean plane through the origin, and
let $e$ be a unit normal to $P$. Choose linearly independent
$a,b\in P$, an interval $J=[\alpha,\beta]$ with $\alpha<\beta$,
a continuous function $R:J\to[0,\infty)$, and a continuously
differentiable function $w:J\to P$. With $p_0\in P$ and $L>0$,
consider the compact set
\begin{equation}\label{eq:radial-sweep}
 \mathcal S=\{p_0+\rho(a+tb)+s(e+w(t)):
          t\in J,\ 0\le\rho\le R(t),\ |s|\le L\}.
\end{equation}
At each $t$, the intercepts form one radial segment in $P$ and the
parallel swept segments have direction $e+w(t)$. The direction varies
with $t$, so this description is not yet a finite cylinder cover.

Put $\kappa=|\det_P(a,b)|$, the area of their parallelogram.
The radial parameterization has area element $\kappa\rho\,d\rho\,dt$.
Projection onto the plane perpendicular to $e+w(t)$ contributes
the factor $(1+\norm{w(t)}^2)^{-1/2}$. These two factors suggest the
weighted area
\begin{equation}\label{eq:weighted-area}
 \mathcal A=\kappa\int_\alpha^\beta
       \frac{R(t)^2}{2\sqrt{1+\norm{w(t)}^2}}\,dt.
\end{equation}
Freezing arbitrary segment directions need not preserve this cost.
The following alignment condition makes the loss vanish.

\begin{theorem}\label{thm:radial}
For the data in \eqref{eq:radial-sweep}, suppose that
\[
 w'(t)=\lambda(t)(a+tb)\qquad(t\in J)
\]
for a continuous real function $\lambda$.
For any finite partition
\(\alpha=t_0<t_1<\cdots<t_N=\beta\), put \(I_i=[t_{i-1},t_i]\)
and choose arbitrary tags \(c_i\in I_i\). There are \(N\) cylinders
covering \(\mathcal S\), with axes parallel to \(e+w(c_i)\) and
nondegenerate compact triangular perpendicular bases. These covers
can be chosen so that their total base area tends to \(\mathcal A\)
as the mesh \(\delta=\max_i|I_i|\) tends to zero, uniformly over the
partitions and tags.

Alternatively, take each base to be the perpendicular projection of
the actual compact sweep with \(t\in I_i\), in the same frozen
direction \(e+w(c_i)\). These \(N\) cylinders cover \(\mathcal S\),
and their total base area also tends to \(\mathcal A\), uniformly
over the partitions and tags.
\end{theorem}

The hypothesis has a direct geometric meaning. To first order, a
change in direction moves the intercept along its existing ray.
After that displacement is absorbed into the radial coordinate, the
remaining transverse displacement is quadratic in the interval length.
A slightly extended triangle contains all these intercepts and adds
only a vanishing total area. The radial cap $R$ may vanish, and
$\lambda$ may change sign or vanish.

The tetrahedral application gives the following concrete consequence.
\begin{corollary}\label{cor:tetrahedron}
Every regular tetrahedron $K\subset\R^3$ of positive edge length has
a finite cylinder cover $B_i+\R u_i$, with unit axes $u_i$ and
nondegenerate compact triangular bases $B_i\subset u_i^\perp$, such that
\[
 \sum_i\area(B_i)<\frac{A_{\min}(K)}2,
 \qquad
 \sum_i\frac{\area(B_i)}{\area(\pi_{u_i^\perp}K)}<\frac12.
\]
\end{corollary}

There is a useful connection with finite approximation for velocity
fields having square-zero derivative. Away from the apex, write
$z=\rho(a+tb)$ with $\rho>0$ and $t\in(\alpha,\beta)$, and define
$W(z)=w(t)$. This is a $C^1$ field in the open angular sector.
The angular differential $Dt$ annihilates $a+tb$, whereas
$DW=w'(t)\otimes Dt$. The alignment hypothesis therefore gives
$(DW)^2=0$. General finite approximation for such fields concerns a
larger class of ruled sets; see the companion article
\cite[Theorem~1.2 and Section~3]{RuledApproximation2026}. Theorem~\ref{thm:radial} proves the
simultaneous prescribed-tag, one-cylinder-per-interval and triangular-base
guarantee directly, including the apex. Merely splitting parallelogram
bases into triangles would not preserve that prescribed count and choice
of axes. No approximation theorem for general velocity fields is used here.

The proof of Theorem~\ref{thm:radial} in
Section~\ref{sec:freezing} constructs the enclosing triangles and
bounds their total excess area. The actual projected sweep contains
the projection of the unperturbed radial sector, which supplies the matching lower
bound. Thus different finite bases have the same limiting cost.

Section~\ref{sec:tetrahedron} proves Corollary~\ref{cor:tetrahedron}
with two overlapping sweeps. A cubic enlargement of their radial caps
exceeds the uniform error in selecting a segment through each point,
while a negative quadratic projection cost survives finite
approximation. Section~\ref{beam:section} develops the alternative
shortened-beam construction: the beam saving and the residual-strip
cost cancel to first order, leaving a strict quadratic improvement.
Its quantitative cost statement separates this construction from the
all-triangular cover already obtained.

The appendices collect supplementary calculations and alternative
finite bases. Appendix~\ref{sec:projections} gives simplex projection
normalizations. Appendix~\ref{sec:cap-choices} compares overlap
margins, Appendix~\ref{sec:finite-footprints} computes explicit finite
enclosures, and Appendix~\ref{sec:coordinate-realizations} transports
the construction between physical and barycentric coordinates.
These calculations retain their own formulas and geometric
realizations; neither main construction requires them.
